\documentclass[aspectratio=43,10pt]{beamer}

% =============================================================================
% 1. PARÁMETROS Y CONFIGURACIÓN DEL PROYECTO (FÁCILMENTE EDITABLE)
% =============================================================================

% Metadatos institucionales
\newcommand{\Universidad}{Universidad Central de Venezuela}
\newcommand{\Facultad}{Facultad de Ciencias}
\newcommand{\Escuela}{Escuela de Computación}
\newcommand{\Materia}{Aprendizaje Profundo}
\newcommand{\Semestre}{Caracas, semestre 2026}
\newcommand{\CodigoMateriaPie}{DL-2026 \quad UCV}

% Datos del equipo docente y estudiantil
\newcommand{\Profesor}{Fernando Crema García}
\newcommand{\EstudianteUno}{Miguel Ciavato}
\newcommand{\EstudianteDos}{Edwyn Guzmán}
\newcommand{\EstudianteTres}{Diego Perdomo}

% Título del proyecto
\newcommand{\TituloPrincipalLuno}{Clasificación de género facial mediante}
\newcommand{\TituloPrincipalLdos}{Vision Transformer y transferencia de aprendizaje}
\newcommand{\TituloPie}{Clasificación de género facial con Vision Transformer}

% Enlaces y repositorio público
\newcommand{\RepoHuggingFaceUrl}{https://huggingface.co/miguelciav/vit-base-utkface-gender}
\newcommand{\RepoHuggingFaceTxt}{miguelciav/vit-base-utkface-gender}

% Rutas de figuras y logotipos
\newcommand{\LogoUCV}{figures/logo_ucv.png}
\newcommand{\LogoIA}{figures/logo_ia.png}
\newcommand{\FigDistribucion}{figures/fig_1_cell_19.png}
\newcommand{\FigMuestras}{figures/fig_2_cell_20.png}
\newcommand{\FigCurvas}{figures/fig_3_cell_50.png}
\newcommand{\FigConfusion}{figures/fig_4_cell_53.png}
\newcommand{\FigErrores}{figures/fig_5_cell_56.png}

% =============================================================================
% 2. PAQUETES Y COMPATIBILIDAD
% =============================================================================
\usepackage[T1]{fontenc}
\usepackage{lmodern}
\usepackage[utf8]{inputenc}
\usepackage[spanish,es-tabla]{babel}
\usepackage{graphicx}
\usepackage{booktabs}
\usepackage{amsmath}
\usepackage{xcolor}
\usepackage{url}

% =============================================================================
% 3. TEMA Y PALETA INSTITUCIONAL UCV
% =============================================================================
\usetheme{Frankfurt}
\usecolortheme{whale}

\definecolor{ucvblue}{RGB}{27, 102, 176}        % Azul institucional medio (barra superior y bloques)
\definecolor{ucvdarkblue}{RGB}{16, 68, 122}      % Azul marino oscuro (centro del pie de página)
\definecolor{ucvlightblue}{RGB}{205, 225, 245}   % Azul claro institucional (laterales del pie de página)
\definecolor{ucvblockbg}{RGB}{242, 246, 252}     % Fondo claro para cajas de contenido
\definecolor{darkgray}{RGB}{35, 40, 45}          % Color de texto legible

\setbeamercolor{palette primary}{bg=ucvdarkblue,fg=white}
\setbeamercolor{palette secondary}{bg=ucvlightblue,fg=ucvdarkblue}
\setbeamercolor{palette tertiary}{bg=ucvblue,fg=white}
\setbeamercolor{section in head/foot}{fg=white,bg=ucvblue}
\setbeamercolor{structure}{fg=ucvblue}
\setbeamercolor{titlelike}{fg=ucvblue}
\setbeamercolor{frametitle}{fg=ucvblue,bg=white}
\setbeamercolor{block title}{bg=ucvblue,fg=white}
\setbeamercolor{block body}{bg=ucvblockbg,fg=darkgray}

% Plantillas para que la navegación superior muestre siempre la sección activa en blanco brillante
\setbeamertemplate{section in head/foot}{\color{white}\textbf{\insertsectionhead}}
\setbeamertemplate{section in head/foot shaded}{\color{white!65!ucvblue}\insertsectionhead}

% Hipervínculos activos y seguros
\hypersetup{
  colorlinks=true,
  linkcolor=white,
  urlcolor=ucvblue,
  citecolor=ucvblue
}

% Desactivar iconos inferiores por defecto de beamer
\setbeamertemplate{navigation symbols}{}

% Tipografía académica de títulos de diapositivas
\setbeamerfont{frametitle}{series=\bfseries,size=\Large,shape=\itshape}

% Pie de página en tres segmentos institucionales
\setbeamertemplate{footline}{%
  \leavevmode%
  \hbox{%
  \begin{beamercolorbox}[wd=.28\paperwidth,ht=2.4ex,dp=1.1ex,left,leftskip=2ex]{palette secondary}%
    \textbf{\CodigoMateriaPie}
  \end{beamercolorbox}%
  \begin{beamercolorbox}[wd=.44\paperwidth,ht=2.4ex,dp=1.1ex,center]{palette primary}%
    \textit{\TituloPie}
  \end{beamercolorbox}%
  \begin{beamercolorbox}[wd=.28\paperwidth,ht=2.4ex,dp=1.1ex,right,rightskip=2ex]{palette secondary}%
    Semestre 2026 \qquad \insertframenumber{} / \inserttotalframenumber
  \end{beamercolorbox}}%
  \vskip0pt%
}

% Macro para láminas de transición de sección con fondo azul sólido
\newcommand{\LaminaTransicion}[1]{%
  {%
    \setbeamercolor{background canvas}{bg=ucvblue}%
    \begin{frame}[plain]%
      \centering\vfill%
      {\Huge \textbf{\color{white} #1}}%
      \vfill%
    \end{frame}%
  }%
}

% Metadatos globales
\title[\TituloPie]{\TituloPrincipalLuno{} \TituloPrincipalLdos}
\author{\EstudianteUno{} \and \EstudianteDos{} \and \EstudianteTres}
\institute{\Escuela, \Facultad, \Universidad}
\date{Semestre 2026}

% =============================================================================
% 4. CONTENIDO DE LA PRESENTACIÓN
% =============================================================================
\begin{document}

% -----------------------------------------------------------------------------
% LÁMINA 1: Portada Institucional
% -----------------------------------------------------------------------------
\begin{frame}[plain]
  \vspace{0.1cm}
  % Fila superior con logotipos y membrete centrado
  \noindent
  \begin{minipage}{0.18\textwidth}
    \raggedright
    \includegraphics[height=1.45cm]{\LogoUCV}
  \end{minipage}%
  \hfill%
  \begin{minipage}{0.60\textwidth}
    \centering
    {\scriptsize\bfseries \Universidad}\\[-0.5pt]
    {\scriptsize \Facultad}\\[-0.5pt]
    {\scriptsize \Escuela}\\[0.5pt]
    {\scriptsize\bfseries\color{ucvblue} \Materia}
  \end{minipage}%
  \hfill%
  \begin{minipage}{0.18\textwidth}
    \raggedleft
    \includegraphics[height=1.45cm]{\LogoIA}
  \end{minipage}

  \vspace{0.45cm}
  % Banner horizontal a pantalla completa
  \noindent\makebox[\linewidth][c]{%
    \begin{beamercolorbox}[wd=\paperwidth,sep=14pt,center]{palette secondary}
      {\Large \textbf{\color{ucvblue}\TituloPrincipalLuno\\[4pt]\TituloPrincipalLdos}}
    \end{beamercolorbox}%
  }

  \vspace{0.45cm}
  % Bloques de integrantes y profesor
  \begin{columns}[t]
    \begin{column}{0.48\textwidth}
      \centering
      {\small \textbf{\color{ucvblue} Integrantes}}\\[7pt]
      {\normalsize \EstudianteUno}\\[3pt]
      {\normalsize \EstudianteDos}\\[3pt]
      {\normalsize \EstudianteTres}
    \end{column}
    \begin{column}{0.48\textwidth}
      \centering
      {\small \textbf{\color{ucvblue} Profesor}}\\[7pt]
      {\normalsize \Profesor}
    \end{column}
  \end{columns}

  \vspace{0.6cm}
  \begin{center}
    {\footnotesize \Semestre}
  \end{center}
\end{frame}

% -----------------------------------------------------------------------------
% LÁMINA 2: Contenido de la Presentación (Agenda de 5 puntos)
% -----------------------------------------------------------------------------
\begin{frame}[t]{Contenido}
  \vspace{0.3cm}
  \begin{enumerate}
    \setlength{\itemsep}{15pt}
    \item \textbf{\large El problema y por qué importa}
    \item \textbf{\large Los datos}
    \item \textbf{\large La tabla de resultados y la discusión}
    \item \textbf{\large Demostración en vivo de la aplicación}
    \item \textbf{\large Qué haríamos distinto y conclusiones}
  \end{enumerate}
\end{frame}

% =============================================================================
% SECCIÓN 1: El problema y por qué importa
% =============================================================================
\section{El problema}

\LaminaTransicion{El problema y por qué importa}

\begin{frame}{El problema y por qué importa}
  \begin{columns}[T]
    \begin{column}{0.48\textwidth}
      \begin{block}{Definición de la tarea}
        \small
        \textbf{Objetivo práctico:} Clasificar automáticamente el género a partir de una fotografía facial para habilitar interfaces adaptativas y analítica demográfica agregada.\\[6pt]
        \textbf{Estrategia de solución:} En lugar de entrenar desde cero, reutilizamos un modelo visual de última generación ya entrenado en millones de imágenes para transferir su conocimiento a rostros humanos.
      \end{block}
    \end{column}
    \hfill
    \begin{column}{0.48\textwidth}
      \begin{block}{Valor para la toma de decisiones}
        \small
        \textbf{Eficiencia operativa:} La transferencia de aprendizaje permite crear soluciones de visión por computador confiables en días en lugar de meses, sin requerir infraestructuras de cómputo millonarias.\\[6pt]
        \textbf{Pregunta clave:} Comprobar si los modelos modernos basados en atención superan a los enfoques tradicionales cuando disponemos de un volumen moderado de datos.
      \end{block}
    \end{column}
  \end{columns}

  \vspace{0.12cm}
  \begin{block}{Consideración ética y uso responsable}
    \footnotesize
    \textbf{Alcance estrictamente técnico:} El sistema detecta patrones visuales superficiales en fotografías para fines estadísticos. No determina ni pretende definir la identidad personal ni la identidad de género de ningún individuo.
  \end{block}
\end{frame}

% =============================================================================
% SECCIÓN 2: Los datos
% =============================================================================
\section{Los datos}

\LaminaTransicion{Los datos}

\begin{frame}{Conjunto de datos UTKFace}
  \begin{columns}[T]
    \begin{column}{0.48\textwidth}
      \begin{block}{Composición y equilibrio del catálogo}
        \scriptsize
        Trabajamos con 18.000 fotografías del conjunto UTKFace seleccionadas con balance equitativo: 52,3\,\% de hombres y 47,7\,\% de mujeres.\\[3pt]
        \textbf{Entrenamiento:} 12.600 imágenes para enseñar al modelo.\\
        \textbf{Validación:} 2.700 imágenes para calibrar el aprendizaje.\\
        \textbf{Evaluación final:} 2.700 imágenes nunca antes vistas para medir la precisión real.
      \end{block}
    \end{column}
    \hfill
    \begin{column}{0.48\textwidth}
      \begin{block}{Control de calidad y preparación}
        \scriptsize
        Estandarizamos todas las imágenes a un tamaño uniforme de $224 \times 224$ píxeles y aplicamos variaciones controladas de rotación y volteo horizontal para que el sistema aprenda a reconocer rostros en condiciones cotidianas.
      \end{block}
    \end{column}
  \end{columns}

  \vspace{0.2cm}
  \begin{columns}[c]
    \begin{column}{0.48\textwidth}
      \centering
      % Gráfico de distribución de clases
      \includegraphics[width=\textwidth,height=3.1cm,keepaspectratio]{\FigDistribucion}
    \end{column}
    \hfill
    \begin{column}{0.48\textwidth}
      \centering
      % Muestra visual de rostros del dataset
      \includegraphics[width=\textwidth,height=3.1cm,keepaspectratio]{\FigMuestras}
    \end{column}
  \end{columns}
\end{frame}

% =============================================================================
% SECCIÓN 3: La tabla de resultados y la discusión
% =============================================================================
\section{Resultados}

\LaminaTransicion{La tabla de resultados y la discusión}

\begin{frame}{Comparación de opciones evaluadas}
  \begin{table}[h]
    \centering
    \resizebox{\textwidth}{!}{%
    \begin{tabular}{lcccc}
      \toprule
      \textbf{Estrategia evaluada} & \textbf{Pesos ajustados} & \textbf{Tiempo} & \textbf{Precisión validación} & \textbf{Precisión prueba} \\
      \midrule
      4.1 Construcción desde cero   & 1.965.890  & 284 s & 69,19\% & 67,04\% \\
      4.2 Modelo base sin adaptar   & 0          & 23 s  & 64,52\% & 65,41\% \\
      4.3 Extracción de rasgos      & 1.538      & 243 s & 94,37\% & 93,56\% \\
      \textbf{4.4 Ajuste fino completo} & \textbf{85.800.194} & \textbf{395 s} & \textbf{95,07\%} & \textbf{94,33\%} \\
      4.5 Ajuste de capas finales   & 14.178.818 & 273 s & 94,78\% & 93,78\% \\
      \bottomrule
    \end{tabular}%
    }
  \end{table}

  \vspace{0.1cm}
  \begin{block}{Conclusiones clave para la toma de decisiones}
    \scriptsize
    \textbf{La opción más precisa:} Ajustar el modelo completo logra la máxima efectividad con 94,33\% de aciertos sobre clientes y usuarios nuevos.\\[3pt]
    \textbf{La opción más eficiente en costos:} Ajustar únicamente la capa final logra 93,56\% de precisión optimizando apenas 1.538 parámetros, ofreciendo el 99\% del rendimiento con un costo de cómputo mínimo.\\[3pt]
    \textbf{Construir desde cero no es viable:} Solo alcanza 67,04\% de acierto, demostrando que reutilizar conocimiento previo es obligatorio para tener éxito.
  \end{block}
\end{frame}

\begin{frame}{Convergencia y matriz de aciertos}
  \begin{columns}[T]
    \begin{column}{0.48\textwidth}
      \centering
      % Curvas de aprendizaje del modelo
      \includegraphics[width=\textwidth,height=3.4cm,keepaspectratio]{\FigCurvas}
      
      \vspace{0.1cm}
      \begin{block}{Comportamiento del aprendizaje}
        \scriptsize
        El modelo aprende de manera fluida y estable. La curva de validación se mantiene firme sin desviarse, asegurando que el sistema generaliza bien y no memoriza fotos.
      \end{block}
    \end{column}
    \hfill
    \begin{column}{0.48\textwidth}
      \centering
      % Matriz de confusión en evaluación final
      \includegraphics[width=\textwidth,height=3.4cm,keepaspectratio]{\FigConfusion}
      
      \vspace{0.1cm}
      \begin{block}{Desempeño equitativo por clase}
        \scriptsize
        Acierta en 95 de cada 100 hombres y en 94 de cada 100 mujeres sobre las 2.700 pruebas independientes, confirmando que la solución no tiene sesgo favorable hacia ningún género.
      \end{block}
    \end{column}
  \end{columns}
\end{frame}

\begin{frame}{Análisis de costo versus beneficio}
  \begin{columns}[T]
    \begin{column}{0.48\textwidth}
      \begin{block}{El intercambio tecnológico}
        \small
        La decisión entre ambas opciones implica un compromiso directo de recursos:\\[6pt]
        \textbf{Costo de optimización:} Ajustar el modelo completo requiere procesar 55.000 veces más parámetros y 62\,\% más tiempo de cómputo.\\[6pt]
        \textbf{Beneficio obtenido:} Una mejora neta de 0,77 puntos porcentuales en la exactitud final.\\[6pt]
        El ajuste de capas finales representa una alternativa intermedia balanceada.
      \end{block}
    \end{column}
    \hfill
    \begin{column}{0.48\textwidth}
      \begin{exampleblock}{Criterio de recomendación de negocio}
        \small
        \textbf{Para soluciones de bajo costo o móviles:} Recomendamos la extracción de rasgos, ya que demanda mínima memoria y se entrena en minutos sin servidores caros.\\[8pt]
        \textbf{Para operaciones de alta exigencia:} Recomendamos el ajuste completo cuando cada décima de precisión marque una diferencia decisiva en el servicio.
      \end{exampleblock}
    \end{column}
  \end{columns}
\end{frame}

% =============================================================================
% SECCIÓN 4: Demostración en vivo de la aplicación
% =============================================================================
\section{Demostración}

\LaminaTransicion{Demostración en vivo de la aplicación}

\begin{frame}{Demostración en vivo de la aplicación}
  \begin{columns}[T]
    \begin{column}{0.48\textwidth}
      \begin{block}{Servicio central en la nube}
        \small
        \textbf{Servidor automatizado:} Provee consultas en tiempo real y verificación de estado operativo.\\[6pt]
        \textbf{Modelo público en Hugging Face:}\\[2pt]
        {\footnotesize \href{\RepoHuggingFaceUrl}{\texttt{\RepoHuggingFaceTxt}}}\\[6pt]
        \textbf{Tiempo de respuesta:} Inferencia optimizada para responder en fracciones de segundo.
      \end{block}
    \end{column}
    \hfill
    \begin{column}{0.48\textwidth}
      \begin{block}{Aplicación interactiva para usuarios}
        \small
        \textbf{Catálogo de pruebas:} Muestras inmediatas de infantes, adultos y adultos mayores.\\[6pt]
        \textbf{Carga de fotografías:} Permite subir cualquier imagen desde el computador.\\[6pt]
        \textbf{Cámara en vivo:} Detección continua con video en directo y visualización de probabilidades.
      \end{block}
    \end{column}
  \end{columns}

  \vspace{0.35cm}
  \begin{center}
    {\large \textbf{\color{ucvblue}Paso a la demostración práctica:}} {\normalsize Abrimos la aplicación web interactiva para ensayar predicciones en directo.}
  \end{center}
\end{frame}

% =============================================================================
% SECCIÓN 5: Qué haríamos distinto y conclusiones
% =============================================================================
\section{Conclusiones}

\LaminaTransicion{Qué haríamos distinto y conclusiones}

\begin{frame}{Análisis de errores y mejoras}
  \begin{columns}[c]
    \begin{column}{0.48\textwidth}
      \centering
      % Muestra de fallos con alta certeza
      \includegraphics[width=\textwidth,height=3.3cm,keepaspectratio]{\FigErrores}\\[4pt]
      {\tiny \textit{Casos donde el sistema falló con alta certeza}}
    \end{column}
    \hfill
    \begin{column}{0.49\textwidth}
      \begin{block}{¿Por qué se equivocó el sistema?}
        \scriptsize
        Al auditar las fallas con mayor certeza, más del 60\% eran fotos mal etiquetadas de origen, como hombres con barba catalogados como mujeres o bebés donde los rasgos aún no están definidos.
      \end{block}

      \vspace{0.08cm}
      \begin{block}{¿Qué haríamos con más datos o cómputo?}
        \scriptsize
        \textbf{Filtrado de errores en datos:} Detectaríamos y corregiríamos fotos mal catalogadas antes de entrenar para que el modelo aprenda sobre información limpia y confiable.\\[3pt]
        \textbf{Modelos a mayor escala y biometría:} Probaríamos modelos más grandes o especializados en rostros como VGGFace2 e InsightFace para medir si justifican su costo de inversión.\\[3pt]
        \textbf{Garantía de equidad demográfica:} Con más volumen de datos, mediríamos el rendimiento por edades y etnias para asegurar un servicio justo y homogéneo en toda la población.
      \end{block}
    \end{column}
  \end{columns}
\end{frame}

\begin{frame}{Conclusiones}
  \begin{block}{Conclusiones ejecutivas}
    \small
    \textbf{Éxito rotundo de la transferencia:} Reutilizar modelos preentrenados permitió superar el 93,5\% de exactitud desde la primera iteración, demostrando que es la vía más rápida y rentable para implementar visión artificial en la organización.\\[8pt]
    \textbf{Solución lista para producción:} Se completó todo el ciclo: validación rigurosa de 5 estrategias, auditoría de fallas, publicación del modelo en la nube y entrega de una aplicación interactiva funcional.
  \end{block}

  \vspace{0.4cm}
  \begin{center}
    {\Large \textbf{\color{ucvblue}¡Muchas gracias por su atención!}}\\[6pt]
    {\small Modelo público disponible en la nube:}\\[2pt]
    {\footnotesize \href{\RepoHuggingFaceUrl}{\texttt{\RepoHuggingFaceTxt}}}\\[8pt]
    {\small \textit{Espacio abierto para preguntas y comentarios del jurado}}
  \end{center}
\end{frame}

\end{document}
